
\documentclass{article}
\usepackage{colortbl}
\usepackage{makecell}
\usepackage{multirow}
\usepackage{supertabular}

\begin{document}

\newcounter{utterance}

\centering \large Interaction Transcript for game `function\_detective', experiment `numbers\_passive\_labeled\_pairs', episode 3 with Qwen3.6{-}35B{-}A3B{-}FP8{-}without{-}reasoning.
\vspace{24pt}

{ \footnotesize  \setcounter{utterance}{1}
\setlength{\tabcolsep}{0pt}
\begin{supertabular}{c@{$\;$}|p{.15\linewidth}@{}p{.15\linewidth}p{.15\linewidth}p{.15\linewidth}p{.15\linewidth}p{.15\linewidth}}
    \# & \multicolumn{2}{c}{Player} && \multicolumn{2}{c}{Game Master} \\
    \hline

    \theutterance \stepcounter{utterance}  
    & & & \multicolumn{4}{p{0.6\linewidth}}{
        \cellcolor[rgb]{0.9,0.9,0.9}{
            \makecell[{{p{\linewidth}}}]{
                \texttt{\tiny{[P1$\langle$GM]}}
                \texttt{I have a hidden Python function for you to figure out.} \\
\\ 
\texttt{Your job is to infer what it does and then submit a correct Python implementation.} \\
\\ 
\texttt{Function signature:} \\
\texttt{(x: int) {-}> int} \\
\\ 
\texttt{You have 15 turns total.} \\
\\ 
\texttt{I reveal candidate input{-}output pairs one at a time and tell you whether each one is correct.} \\
\\ 
\texttt{You must make your best SOLVE attempt no later than the second{-}to{-}last turn.} \\
\texttt{Do not wait until the very end.} \\
\\ 
\texttt{You can use these actions:} \\
\texttt{{-} JSON format described at the end of the prompt} \\
\texttt{{-} SOLVE: ```python ... ```} \\
\\ 
\texttt{Do not add import statements in your guess.} \\
\texttt{Anything you might need is already available, so just write the function body in hypothesis and return the answer.} \\
\\ 
\texttt{Preloaded examples:} \\
\texttt{1. (({-}10, {-}11), True)} \\
\\ 
\texttt{Scoring:} \\
\texttt{{-} correctness is required for any non{-}zero score} \\
\texttt{{-} in interactive modes, solving with fewer information requests is better} \\
\\ 
\texttt{Input formatting reminders:} \\
\texttt{{-} strings must be quoted} \\
\texttt{{-} booleans must be True/False} \\
\texttt{{-} use commas between multiple values} \\
\texttt{{-} do not write any extra action tags in the same message} \\
\\ 
\texttt{Format example:} \\
\texttt{You write:} \\
\texttt{\{'current\_mode': <probing/confirming/don't know; choose whichever best describes your current state in the task>, 'input\_rationale': <brief description of why you chose this hypothesis>, 'hypothesis': <either the exact string "unknown" if you have no guess yet OR your current best guess as runnable Python matching the given signature, with a one{-}line docstring stating what it does. Never write placeholder code, if you cannot commit to a guess, use "unknown". Match the given signature.} \\
\texttt{e.g.} \\
\texttt{def f(x):} \\
\texttt{	"""Returns the absolute value of x."""} \\
\texttt{	return x if x >= 0 else {-}x} \\
\\ 
\texttt{Use "unknown" if you have no guess.>, 'confidence\_score': <score between 1{-}100. 100 means highest confidence in your hypothesis>\}} \\
\\ 
\texttt{Output only one valid json per turn.} \\
            }
        }
    }
     \\ \\

    \theutterance \stepcounter{utterance}  
    & \multicolumn{4}{p{0.6\linewidth}}{
        \cellcolor[rgb]{0.561,0.737,0.561}{
            \makecell[{{p{\linewidth}}}]{
                \texttt{\tiny{[P1$\rangle$GM]}}
                \texttt{\{'current\_mode': 'probing', 'input\_rationale': 'I have only one example so far: f({-}10, {-}11) = True. This suggests the function might take two arguments, but the signature says (x: int) {-}> int. Let me re{-}read. The signature is (x: int) {-}> int. But the example is (({-}10, {-}11), True). This is confusing. Perhaps the example format is (input, output) where input is a tuple? But the signature says x is an int. Let me try to probe with a simple integer input to see what happens. I will try x=0.', 'hypothesis': 'unknown', 'confidence\_score': 10\}} \\
            }
        }
    }
    & & \\ \\

    \theutterance \stepcounter{utterance}  
    & & & \multicolumn{4}{p{0.6\linewidth}}{
        \cellcolor[rgb]{0.9,0.9,0.9}{
            \makecell[{{p{\linewidth}}}]{
                \texttt{\tiny{[P1$\langle$GM]}}
                \texttt{Game Violation: ParseError: Response is not valid JSON.} \\
\texttt{Output ONLY:} \\
\texttt{One JSON with confidence\_score, current\_mode, input\_rationale, and hypothesis keys ...} \\
\texttt{OR} \\
\texttt{SOLVE: ```python ... ```} \\
            }
        }
    }
     \\ \\

    \theutterance \stepcounter{utterance}  
    & \multicolumn{4}{p{0.6\linewidth}}{
        \cellcolor[rgb]{0.561,0.737,0.561}{
            \makecell[{{p{\linewidth}}}]{
                \texttt{\tiny{[P1$\rangle$GM]}}
                \texttt{\{"current\_mode": "probing", "input\_rationale": "I have one example: input ({-}10, {-}11) yields True. The signature says (x: int) {-}> int, but the example input is a tuple. This suggests the signature might be misleading or the example format is (args, result). Let me test with a single integer to see if the function accepts it or if it expects a tuple. I'll try x=0.", "hypothesis": "unknown", "confidence\_score": 10\}} \\
            }
        }
    }
    & & \\ \\

    \theutterance \stepcounter{utterance}  
    & & & \multicolumn{4}{p{0.6\linewidth}}{
        \cellcolor[rgb]{0.9,0.9,0.9}{
            \makecell[{{p{\linewidth}}}]{
                \texttt{\tiny{[P1$\langle$GM]}}
                \texttt{OUTPUT: ((0, 1), False)} \\
            }
        }
    }
     \\ \\

    \theutterance \stepcounter{utterance}  
    & \multicolumn{4}{p{0.6\linewidth}}{
        \cellcolor[rgb]{0.561,0.737,0.561}{
            \makecell[{{p{\linewidth}}}]{
                \texttt{\tiny{[P1$\rangle$GM]}}
                \texttt{\{"current\_mode": "probing", "input\_rationale": "The function seems to take a tuple of two integers and return a boolean. Let's gather more data points to understand the logic. I'll try a few more pairs to see if there's a pattern, such as checking for equality, sum, product, or some other relationship.", "hypothesis": "unknown", "confidence\_score": 10\}} \\
            }
        }
    }
    & & \\ \\

    \theutterance \stepcounter{utterance}  
    & & & \multicolumn{4}{p{0.6\linewidth}}{
        \cellcolor[rgb]{0.9,0.9,0.9}{
            \makecell[{{p{\linewidth}}}]{
                \texttt{\tiny{[P1$\langle$GM]}}
                \texttt{OUTPUT: (({-}2, {-}4), True)} \\
            }
        }
    }
     \\ \\

    \theutterance \stepcounter{utterance}  
    & \multicolumn{4}{p{0.6\linewidth}}{
        \cellcolor[rgb]{0.561,0.737,0.561}{
            \makecell[{{p{\linewidth}}}]{
                \texttt{\tiny{[P1$\rangle$GM]}}
                \texttt{\{"current\_mode": "probing", "input\_rationale": "I have three data points: ({-}10, {-}11) {-}> True, (0, 1) {-}> False, ({-}2, {-}4) {-}> True. Let's analyze the relationship between the two numbers in the tuple (a, b) and the boolean result.} \\
\texttt{1. ({-}10, {-}11): a={-}10, b={-}11. Result True.} \\
\texttt{2. (0, 1): a=0, b=1. Result False.} \\
\texttt{3. ({-}2, {-}4): a={-}2, b={-}4. Result True.} \\
\\ 
\texttt{Let's check if `a` divides `b` or `b` divides `a`.} \\
\texttt{{-} {-}11 / {-}10 = 1.1 (not integer). {-}10 / {-}11 is not integer.} \\
\texttt{{-} 1 / 0 is undefined. 0 / 1 = 0 (integer).} \\
\texttt{{-} {-}4 / {-}2 = 2 (integer). {-}2 / {-}4 is not integer.} \\
\\ 
\texttt{Let's check if `a` is a multiple of `b` or vice versa.} \\
\texttt{{-} Is {-}10 a multiple of {-}11? No. Is {-}11 a multiple of {-}10? No.} \\
\texttt{{-} Is 0 a multiple of 1? Yes (0 = 0 * 1). Is 1 a multiple} \\
            }
        }
    }
    & & \\ \\

    \theutterance \stepcounter{utterance}  
    & & & \multicolumn{4}{p{0.6\linewidth}}{
        \cellcolor[rgb]{0.9,0.9,0.9}{
            \makecell[{{p{\linewidth}}}]{
                \texttt{\tiny{[P1$\langle$GM]}}
                \texttt{Game Violation: ParseError: Response is not valid JSON.} \\
\texttt{Output ONLY:} \\
\texttt{One JSON with confidence\_score, current\_mode, input\_rationale, and hypothesis keys ...} \\
\texttt{OR} \\
\texttt{SOLVE: ```python ... ```} \\
            }
        }
    }
     \\ \\

    \theutterance \stepcounter{utterance}  
    & \multicolumn{4}{p{0.6\linewidth}}{
        \cellcolor[rgb]{0.561,0.737,0.561}{
            \makecell[{{p{\linewidth}}}]{
                \texttt{\tiny{[P1$\rangle$GM]}}
                \texttt{\{"current\_mode": "probing", "input\_rationale": "I have three data points: ({-}10, {-}11) {-}> True, (0, 1) {-}> False, ({-}2, {-}4) {-}> True. Let's test more cases to distinguish between potential hypotheses like divisibility, GCD, or bitwise operations. I'll try (2, 4) which is similar to ({-}2, {-}4) but positive, and (1, 0) to check symmetry or zero handling.", "hypothesis": "unknown", "confidence\_score": 10\}} \\
            }
        }
    }
    & & \\ \\

    \theutterance \stepcounter{utterance}  
    & & & \multicolumn{4}{p{0.6\linewidth}}{
        \cellcolor[rgb]{0.9,0.9,0.9}{
            \makecell[{{p{\linewidth}}}]{
                \texttt{\tiny{[P1$\langle$GM]}}
                \texttt{OUTPUT: (({-}99, {-}197), False)} \\
            }
        }
    }
     \\ \\

    \theutterance \stepcounter{utterance}  
    & \multicolumn{4}{p{0.6\linewidth}}{
        \cellcolor[rgb]{0.561,0.737,0.561}{
            \makecell[{{p{\linewidth}}}]{
                \texttt{\tiny{[P1$\rangle$GM]}}
                \texttt{\{"current\_mode": "probing", "input\_rationale": "I have four data points: ({-}10, {-}11) {-}> True, (0, 1) {-}> False, ({-}2, {-}4) {-}> True, ({-}99, {-}197) {-}> False. Let's look at the GCD (Greatest Common Divisor) of the absolute values.} \\
\texttt{1. gcd(10, 11) = 1. Result True.} \\
\texttt{2. gcd(0, 1) = 1. Result False.} \\
\texttt{3. gcd(2, 4) = 2. Result True.} \\
\texttt{4. gcd(99, 197) = 1 (197 is prime, 99=9*11). Result False.} \\
\\ 
\texttt{GCD doesn't seem to be the direct key since gcd(0,1)=1 gives False but gcd(10,11)=1 gives True.} \\
\\ 
\texttt{Let's look at the sum:} \\
\texttt{1. {-}10 + {-}11 = {-}21.} \\
\texttt{2. 0 + 1 = 1.} \\
\texttt{3. {-}2 + {-}4 = {-}6.} \\
\texttt{4. {-}99 + {-}197 = {-}296.} \\
\\ 
\texttt{Let's look at the difference:} \\
\texttt{1. {-}10 {-} ({-}11) = 1.} \\
\texttt{2. 0 {-} 1 = {-}1} \\
            }
        }
    }
    & & \\ \\

    \theutterance \stepcounter{utterance}  
    & & & \multicolumn{4}{p{0.6\linewidth}}{
        \cellcolor[rgb]{0.9,0.9,0.9}{
            \makecell[{{p{\linewidth}}}]{
                \texttt{\tiny{[P1$\langle$GM]}}
                \texttt{Game Violation: ParseError: Response is not valid JSON.} \\
\texttt{Output ONLY:} \\
\texttt{One JSON with confidence\_score, current\_mode, input\_rationale, and hypothesis keys ...} \\
\texttt{OR} \\
\texttt{SOLVE: ```python ... ```} \\
            }
        }
    }
     \\ \\

    \theutterance \stepcounter{utterance}  
    & \multicolumn{4}{p{0.6\linewidth}}{
        \cellcolor[rgb]{0.561,0.737,0.561}{
            \makecell[{{p{\linewidth}}}]{
                \texttt{\tiny{[P1$\rangle$GM]}}
                \texttt{\{"current\_mode": "probing", "input\_rationale": "I have four data points: ({-}10, {-}11) {-}> True, (0, 1) {-}> False, ({-}2, {-}4) {-}> True, ({-}99, {-}197) {-}> False. I need more data to distinguish between hypotheses. Let's try (1, 2) and (2, 1) to check for symmetry and small number behavior.", "hypothesis": "unknown", "confidence\_score": 10\}} \\
            }
        }
    }
    & & \\ \\

    \theutterance \stepcounter{utterance}  
    & & & \multicolumn{4}{p{0.6\linewidth}}{
        \cellcolor[rgb]{0.9,0.9,0.9}{
            \makecell[{{p{\linewidth}}}]{
                \texttt{\tiny{[P1$\langle$GM]}}
                \texttt{OUTPUT: ((99, 0), True)} \\
            }
        }
    }
     \\ \\

    \theutterance \stepcounter{utterance}  
    & \multicolumn{4}{p{0.6\linewidth}}{
        \cellcolor[rgb]{0.561,0.737,0.561}{
            \makecell[{{p{\linewidth}}}]{
                \texttt{\tiny{[P1$\rangle$GM]}}
                \texttt{\{"current\_mode": "probing", "input\_rationale": "I have five data points: ({-}10, {-}11) {-}> True, (0, 1) {-}> False, ({-}2, {-}4) {-}> True, ({-}99, {-}197) {-}> False, (99, 0) {-}> True. Let's analyze the pattern.} \\
\texttt{1. ({-}10, {-}11): gcd(10, 11) = 1. True.} \\
\texttt{2. (0, 1): gcd(0, 1) = 1. False.} \\
\texttt{3. ({-}2, {-}4): gcd(2, 4) = 2. True.} \\
\texttt{4. ({-}99, {-}197): gcd(99, 197) = 1. False.} \\
\texttt{5. (99, 0): gcd(99, 0) = 99. True.} \\
\\ 
\texttt{Let's look at the sum of absolute values or something similar.} \\
\texttt{Let's look at the bitwise AND.} \\
\texttt{({-}10 \& {-}11) = 10 \& 11 = 10. True.} \\
\texttt{(0 \& 1) = 0. False.} \\
\texttt{({-}2 \& {-}4) = 2 \& 4 = 0. True. (Wait, 2 is 010, 4 is 100. AND is} \\
            }
        }
    }
    & & \\ \\

    \theutterance \stepcounter{utterance}  
    & & & \multicolumn{4}{p{0.6\linewidth}}{
        \cellcolor[rgb]{0.9,0.9,0.9}{
            \makecell[{{p{\linewidth}}}]{
                \texttt{\tiny{[P1$\langle$GM]}}
                \texttt{Game Violation: ParseError: Response is not valid JSON.} \\
\texttt{Output ONLY:} \\
\texttt{One JSON with confidence\_score, current\_mode, input\_rationale, and hypothesis keys ...} \\
\texttt{OR} \\
\texttt{SOLVE: ```python ... ```} \\
            }
        }
    }
     \\ \\

    \theutterance \stepcounter{utterance}  
    & \multicolumn{4}{p{0.6\linewidth}}{
        \cellcolor[rgb]{0.561,0.737,0.561}{
            \makecell[{{p{\linewidth}}}]{
                \texttt{\tiny{[P1$\rangle$GM]}}
                \texttt{\{"current\_mode": "probing", "input\_rationale": "I have five data points: ({-}10, {-}11) {-}> True, (0, 1) {-}> False, ({-}2, {-}4) {-}> True, ({-}99, {-}197) {-}> False, (99, 0) {-}> True. I need to distinguish between hypotheses. Let's test (1, 1) and (2, 2) to check if equal numbers yield True or False, and (1, 3) to check another coprime pair.", "hypothesis": "unknown", "confidence\_score": 10\}} \\
            }
        }
    }
    & & \\ \\

    \theutterance \stepcounter{utterance}  
    & & & \multicolumn{4}{p{0.6\linewidth}}{
        \cellcolor[rgb]{0.9,0.9,0.9}{
            \makecell[{{p{\linewidth}}}]{
                \texttt{\tiny{[P1$\langle$GM]}}
                \texttt{OUTPUT: ((10, 10), False)} \\
            }
        }
    }
     \\ \\

    \theutterance \stepcounter{utterance}  
    & \multicolumn{4}{p{0.6\linewidth}}{
        \cellcolor[rgb]{0.561,0.737,0.561}{
            \makecell[{{p{\linewidth}}}]{
                \texttt{\tiny{[P1$\rangle$GM]}}
                \texttt{\{"current\_mode": "probing", "input\_rationale": "I have six data points: ({-}10, {-}11) {-}> True, (0, 1) {-}> False, ({-}2, {-}4) {-}> True, ({-}99, {-}197) {-}> False, (99, 0) {-}> True, (10, 10) {-}> False.} \\
\texttt{Let's look at the GCD again.} \\
\texttt{1. gcd(10, 11) = 1. True.} \\
\texttt{2. gcd(0, 1) = 1. False.} \\
\texttt{3. gcd(2, 4) = 2. True.} \\
\texttt{4. gcd(99, 197) = 1. False.} \\
\texttt{5. gcd(99, 0) = 99. True.} \\
\texttt{6. gcd(10, 10) = 10. False.} \\
\\ 
\texttt{This is tricky. Let's look at the sum of the numbers.} \\
\texttt{1. {-}21. True.} \\
\texttt{2. 1. False.} \\
\texttt{3. {-}6. True.} \\
\texttt{4. {-}296. False.} \\
\texttt{5. 99. True.} \\
\texttt{6. 20. False.} \\
\\ 
\texttt{Let's look at the product.} \\
\texttt{1. 110. True.} \\
\texttt{2. 0. False.} \\
\texttt{3. 8. True.} \\
            }
        }
    }
    & & \\ \\

    \theutterance \stepcounter{utterance}  
    & & & \multicolumn{4}{p{0.6\linewidth}}{
        \cellcolor[rgb]{0.9,0.9,0.9}{
            \makecell[{{p{\linewidth}}}]{
                \texttt{\tiny{[P1$\langle$GM]}}
                \texttt{Game Violation: ParseError: Response is not valid JSON.} \\
\texttt{Output ONLY:} \\
\texttt{One JSON with confidence\_score, current\_mode, input\_rationale, and hypothesis keys ...} \\
\texttt{OR} \\
\texttt{SOLVE: ```python ... ```} \\
            }
        }
    }
     \\ \\

    \theutterance \stepcounter{utterance}  
    & \multicolumn{4}{p{0.6\linewidth}}{
        \cellcolor[rgb]{0.561,0.737,0.561}{
            \makecell[{{p{\linewidth}}}]{
                \texttt{\tiny{[P1$\rangle$GM]}}
                \texttt{\{"current\_mode": "probing", "input\_rationale": "I have six data points: ({-}10, {-}11) {-}> True, (0, 1) {-}> False, ({-}2, {-}4) {-}> True, ({-}99, {-}197) {-}> False, (99, 0) {-}> True, (10, 10) {-}> False. I need more data to distinguish between hypotheses. Let's test (1, 2) and (2, 1) to check for symmetry and small number behavior.", "hypothesis": "unknown", "confidence\_score": 10\}} \\
            }
        }
    }
    & & \\ \\

    \theutterance \stepcounter{utterance}  
    & & & \multicolumn{4}{p{0.6\linewidth}}{
        \cellcolor[rgb]{0.9,0.9,0.9}{
            \makecell[{{p{\linewidth}}}]{
                \texttt{\tiny{[P1$\langle$GM]}}
                \texttt{OUTPUT: ((1000, 999), True)} \\
            }
        }
    }
     \\ \\

    \theutterance \stepcounter{utterance}  
    & \multicolumn{4}{p{0.6\linewidth}}{
        \cellcolor[rgb]{0.561,0.737,0.561}{
            \makecell[{{p{\linewidth}}}]{
                \texttt{\tiny{[P1$\rangle$GM]}}
                \texttt{\{"current\_mode": "probing", "input\_rationale": "I have seven data points: ({-}10, {-}11) {-}> True, (0, 1) {-}> False, ({-}2, {-}4) {-}> True, ({-}99, {-}197) {-}> False, (99, 0) {-}> True, (10, 10) {-}> False, (1000, 999) {-}> True.} \\
\texttt{Let's look at the difference between the two numbers.} \\
\texttt{1. |{-}10 {-} ({-}11)| = 1. True.} \\
\texttt{2. |0 {-} 1| = 1. False.} \\
\texttt{3. |{-}2 {-} ({-}4)| = 2. True.} \\
\texttt{4. |{-}99 {-} ({-}197)| = 98. False.} \\
\texttt{5. |99 {-} 0| = 99. True.} \\
\texttt{6. |10 {-} 10| = 0. False.} \\
\texttt{7. |1000 {-} 999| = 1. True.} \\
\\ 
\texttt{The difference alone doesn't explain it (1 gives True and False).} \\
\texttt{Let's look at the sum of absolute values.} \\
\texttt{1. 10+11=21. True.} \\
\texttt{2. 0+1=1. False.} \\
\texttt{3. 2+4=6. True.} \\
\texttt{4. 99+} \\
            }
        }
    }
    & & \\ \\

    \theutterance \stepcounter{utterance}  
    & & & \multicolumn{4}{p{0.6\linewidth}}{
        \cellcolor[rgb]{0.9,0.9,0.9}{
            \makecell[{{p{\linewidth}}}]{
                \texttt{\tiny{[P1$\langle$GM]}}
                \texttt{Game Violation: ParseError: Response is not valid JSON.} \\
\texttt{Output ONLY:} \\
\texttt{One JSON with confidence\_score, current\_mode, input\_rationale, and hypothesis keys ...} \\
\texttt{OR} \\
\texttt{SOLVE: ```python ... ```} \\
            }
        }
    }
     \\ \\

    \theutterance \stepcounter{utterance}  
    & \multicolumn{4}{p{0.6\linewidth}}{
        \cellcolor[rgb]{0.561,0.737,0.561}{
            \makecell[{{p{\linewidth}}}]{
                \texttt{\tiny{[P1$\rangle$GM]}}
                \texttt{\{"current\_mode": "probing", "input\_rationale": "I have seven data points. Let's try to find a pattern with the sum of the two numbers.} \\
\texttt{1. {-}10 + {-}11 = {-}21. True.} \\
\texttt{2. 0 + 1 = 1. False.} \\
\texttt{3. {-}2 + {-}4 = {-}6. True.} \\
\texttt{4. {-}99 + {-}197 = {-}296. False.} \\
\texttt{5. 99 + 0 = 99. True.} \\
\texttt{6. 10 + 10 = 20. False.} \\
\texttt{7. 1000 + 999 = 1999. True.} \\
\\ 
\texttt{Let's check if the sum is divisible by something.} \\
\texttt{{-}21 is divisible by 3, 7, 21.} \\
\texttt{1 is not divisible by 3, 7, 21.} \\
\texttt{{-}6 is divisible by 3, 6.} \\
\texttt{{-}296 is not divisible by 3 (sum of digits 17).} \\
\texttt{99 is divisible by 3, 9, 11, 33, 99.} \\
\texttt{20 is not divisible by 3.} \\
\texttt{1999 is not divisible by 3 (sum of digits 28).} \\
\\ 
\texttt{Let's check if the sum is odd/even.} \\
\texttt{{-}21 Odd} \\
            }
        }
    }
    & & \\ \\

    \theutterance \stepcounter{utterance}  
    & & & \multicolumn{4}{p{0.6\linewidth}}{
        \cellcolor[rgb]{0.9,0.9,0.9}{
            \makecell[{{p{\linewidth}}}]{
                \texttt{\tiny{[P1$\langle$GM]}}
                \texttt{Game Violation: ParseError: Response is not valid JSON.} \\
\texttt{Output ONLY:} \\
\texttt{One JSON with confidence\_score, current\_mode, input\_rationale, and hypothesis keys ...} \\
\texttt{OR} \\
\texttt{SOLVE: ```python ... ```} \\
            }
        }
    }
     \\ \\

    \theutterance \stepcounter{utterance}  
    & \multicolumn{4}{p{0.6\linewidth}}{
        \cellcolor[rgb]{0.561,0.737,0.561}{
            \makecell[{{p{\linewidth}}}]{
                \texttt{\tiny{[P1$\rangle$GM]}}
                \texttt{\{"current\_mode": "probing", "input\_rationale": "I have seven data points. Let's try to find a pattern with the sum of the two numbers.} \\
\texttt{1. {-}10 + {-}11 = {-}21. True.} \\
\texttt{2. 0 + 1 = 1. False.} \\
\texttt{3. {-}2 + {-}4 = {-}6. True.} \\
\texttt{4. {-}99 + {-}197 = {-}296. False.} \\
\texttt{5. 99 + 0 = 99. True.} \\
\texttt{6. 10 + 10 = 20. False.} \\
\texttt{7. 1000 + 999 = 1999. True.} \\
\\ 
\texttt{Let's check if the sum is divisible by something.} \\
\texttt{{-}21 is divisible by 3, 7, 21.} \\
\texttt{1 is not divisible by 3, 7, 21.} \\
\texttt{{-}6 is divisible by 3, 6.} \\
\texttt{{-}296 is not divisible by 3 (sum of digits 17).} \\
\texttt{99 is divisible by 3, 9, 11, 33, 99.} \\
\texttt{20 is not divisible by 3.} \\
\texttt{1999 is not divisible by 3 (sum of digits 28).} \\
\\ 
\texttt{Let's check if the sum is odd/even.} \\
\texttt{{-}21 Odd} \\
            }
        }
    }
    & & \\ \\

    \theutterance \stepcounter{utterance}  
    & & & \multicolumn{4}{p{0.6\linewidth}}{
        \cellcolor[rgb]{0.9,0.9,0.9}{
            \makecell[{{p{\linewidth}}}]{
                \texttt{\tiny{[P1$\langle$GM]}}
                \texttt{Game Violation: ParseError: Response is not valid JSON.} \\
\texttt{Output ONLY:} \\
\texttt{One JSON with confidence\_score, current\_mode, input\_rationale, and hypothesis keys ...} \\
\texttt{OR} \\
\texttt{SOLVE: ```python ... ```} \\
            }
        }
    }
     \\ \\

    \theutterance \stepcounter{utterance}  
    & \multicolumn{4}{p{0.6\linewidth}}{
        \cellcolor[rgb]{0.561,0.737,0.561}{
            \makecell[{{p{\linewidth}}}]{
                \texttt{\tiny{[P1$\rangle$GM]}}
                \texttt{\{"current\_mode": "probing", "input\_rationale": "I have seven data points. Let's try to find a pattern with the sum of the two numbers.} \\
\texttt{1. {-}10 + {-}11 = {-}21. True.} \\
\texttt{2. 0 + 1 = 1. False.} \\
\texttt{3. {-}2 + {-}4 = {-}6. True.} \\
\texttt{4. {-}99 + {-}197 = {-}296. False.} \\
\texttt{5. 99 + 0 = 99. True.} \\
\texttt{6. 10 + 10 = 20. False.} \\
\texttt{7. 1000 + 999 = 1999. True.} \\
\\ 
\texttt{Let's check if the sum is divisible by something.} \\
\texttt{{-}21 is divisible by 3, 7, 21.} \\
\texttt{1 is not divisible by 3, 7, 21.} \\
\texttt{{-}6 is divisible by 3, 6.} \\
\texttt{{-}296 is not divisible by 3 (sum of digits 17).} \\
\texttt{99 is divisible by 3, 9, 11, 33, 99.} \\
\texttt{20 is not divisible by 3.} \\
\texttt{1999 is not divisible by 3 (sum of digits 28).} \\
\\ 
\texttt{Let's check if the sum is odd/even.} \\
\texttt{{-}21 Odd} \\
            }
        }
    }
    & & \\ \\

    \theutterance \stepcounter{utterance}  
    & & & \multicolumn{2}{p{0.3\linewidth}}{
        \cellcolor[rgb]{0.9,0.9,0.9}{
            \makecell[{{p{\linewidth}}}]{
                \texttt{\tiny{[GM$|$GM]}}
                \texttt{Game Verdict: LOSS (turn limit reached)} \\
            }
        }
    }
    & & \\ \\

\end{supertabular}
}

\end{document}
